\subsection{The residual Bernoulli-unit input in the actual data}
\label{subsec:sr-residual-bernoulli-input}
The arithmetic data constructors use the following residual
vanishing, called \(\mathrm{CorX}\) in the archive.
We prove its explicit Bernoulli estimate and its precise
unramified-cocycle consequence.

\begin{srlemma}\label{lem:sr-dyadic-bernoulli-unit}
Let \(\ell\) be odd, \(d\ge1\),
\(Q=\ell^d\equiv1\pmod4\), and \(e=v_2(Q-1)\).
Let \(f\) be odd, let \(\phi\) be primitive of conductor
\(f\) and order dividing \(Q-1\), and let \(\xi\) be
primitive of conductor \(2^n\).
Suppose
\[
                 f4^e\le2^n,\qquad
                         \phi(-1)\xi(-1)=-1.
 \tag{BX1}
\]
Then \(B_{1,\phi\xi}\) is an \(\ell\)-adic unit.
The statement permits \(\ell\mid f\).
\end{srlemma}
\begin{proof}
Work in a finite splitting extension, with integer ring
\(\mathcal O\) and residue field embedded in an algebraic
closure of \(\mathbb F_\ell\).
Put \(N_2=2^n\), \(m=n-e\).
Equation (BX1) gives
\[
               e\ge2,\qquad m\ge e,\qquad
                         s=f2^e\le2^m .
 \tag{BX2}
\]
Extend both characters by zero on nonunits, as usual.

We first construct an integral function \(g\) on
\(\mathbb Z/N_2\).
Put
\[
 A(u)=\sum_{\substack{0\le b<fN_2\\b\equiv u\ (N_2)}}\phi(b)b,
 \qquad C=\frac1{N_2}\sum_{b=0}^{f-1}\phi(b),\qquad
                  g(u)=\frac{A(u)-A(0)}{fN_2}.
 \tag{BX3}
\]
Although the displayed denominator contains \(f\), every
\(g(u)\) is integral.
Indeed shifting the range of \(b\) by \(f\), using
\(\phi(b+f)=\phi(b)\), gives
\[
              g(u)-g(u-f)
                       =C-\mathbf1_{u_0<f}\phi(u_0),
 \tag{BX4}
\]
where \(u_0\in[0,N_2)\) is the representative.
The constant \(C\) is integral because \(N_2\) is an
\(\ell\)-adic unit.
Addition of \(f\) cycles through every residue modulo
\(N_2\), and \(g(0)=0\).
Thus (BX4) proves integrality at all residues without
dividing by \(f\).
It also shows that the reduction \(G(u)=\bar g(u)\)
lies in \(\mathbb F_Q\), since every \(\bar\phi(b)\)
does and \(N_2^{-1}\in\mathbb F_\ell\).
Telescoping (BX4), whenever \(ft\le N_2\), gives
\[
        G(u)-G(u-ft)
               =t\bar C-\mathbf1_{u_0<ft}\bar\phi(u_0).
 \tag{BX5}
\]
The finite interval contains at most one boundary crossing;
its \(\phi\)-value is unchanged by the multiples of \(f\).
This also verifies (BX5) directly at the boundary.

Since \(\phi\xi\) is nontrivial and odd, its character sum
is zero, and the elementary Bernoulli formula gives
\[
 B_{1,\phi\xi}
   =\frac1{fN_2}\sum_{b=0}^{fN_2-1}\phi(b)\xi(b)b
   =\sum_{u\bmod N_2}\xi(u)g(u).
 \tag{BX6}
\]
The last equality follows by grouping the first sum by
\(u\bmod N_2\) and using \(\sum_u\xi(u)=0\).
In particular the Bernoulli number is integral.

Write \(X=\bar\xi\).
Primitivity gives
\[
             X(1+N_2/2)=-1,
 \tag{BX7}
\]
and its restriction to \(1+4\mathbb Z/2^n\) is faithful.
Suppose for contradiction the reduction of (BX6) is zero.
For every odd \(y\bmod N_2\), changing variables gives
\(\sum_uX(u)G(uy)=0\).
Raise this identity to the powers \(Q^j\).
The coefficients \(G(uy)\) are fixed by that Frobenius.
The powers of \(Q\) modulo \(2^n\) run through
\(1+2^e\mathbb Z/2^n\):
\(v_2(Q^{2^j}-1)=e+j\), by the elementary
difference-of-squares induction.
Averaging these identities with the geometric-sum
identity therefore gives
\[
 H(y):=\sum_{\substack{x\bmod N_2\\X(x)^{2^e}=1}}
                                      X(x)G(xy)=0
                         \qquad(y\ \text{odd}).
 \tag{BX8}
\]
More explicitly average the exponents \(1+2^ev\),
\(0\le v<N_2\).
The inner sum is \(N_2X(x)\) if \(X(x)^{2^e}=1\)
and zero otherwise; \(N_2\) is invertible modulo \(\ell\).

Let \(\mathcal C\) be the indexing subgroup in (BX8).
The decomposition
\((\mathbb Z/2^n)^\times=\{\pm1\}\times\langle5\rangle\)
and \(v_2(5^{2^j}-1)=j+2\) give
\[
       \mathcal C
       =\{x:x\equiv1\text{ or }-1\pmod{2^m}\}.
 \tag{BX9}
\]
For example \(X(5)\) has exact order \(2^{n-2}\);
writing \(x=(-1)^\delta5^j\), the condition is
\(2^{m-2}\mid j\), exactly (BX9).

Apply (BX5) with \(t=2^e\), so \(ft=s\), and subtract
\(H(1-s)\) from \(H(1)\).
For \(x\equiv1\pmod{2^m}\),
\(x(1-s)=x-s\pmod{N_2}\).
Among these \(x\), the interval \(0\le x_0<s\)
contains only \(x=1\).
For \(x\equiv-1\pmod{2^m}\),
\(x(1-s)=x+s\pmod{N_2}\).
The corresponding interval after addition of \(s\)
contains only \(x=-1\), with representative \(s-1\).
The two boundary terms are respectively
\[
                    -X(1)\bar\phi(1)=-1,\qquad
                   X(-1)\bar\phi(s-1)=-1 .
\]
The second equality uses \(f\mid s\) and the oddness
in (BX1).
Thus
\[
 H(1)-H(1-s)
  =2^e\bar C\sum_{x\in\mathcal C}\varepsilon(x)X(x)-2,
 \tag{BX10}
\]
where \(\varepsilon(x)=1\) or \(-1\) according to its
class in (BX9).
The sum in (BX10) is zero.
Indeed multiplication by \(1+N_2/2\) preserves
\(\mathcal C\) and \(\varepsilon(x)\), and multiplies
\(X(x)\) by \(-1\), by (BX7).
Since \(2\ne0\), the sum equals its negative only if zero.
Equations (BX8) and (BX10) now give \(0=-2\),
a contradiction.
The integral value in (BX6) has nonzero reduction,
which proves it is an \(\ell\)-adic unit.
\end{proof}

\begin{srlemma}\label{lem:sr-bernoulli-unit-residual-unramified}
Let \(\chi\ne1\) be a finite Dirichlet character of order
prime to odd \(\ell\), with values in a splitting DVR
\(\mathcal O\).
Let \(\gamma=\bar\chi:G_{\mathbb Q}\to k^\times\),
where \(k\) is any field of characteristic \(\ell\).
Assume
\[
          \gamma\ne\bar\epsilon_\ell,\qquad
                         B_{1,\chi^{-1}}\in\mathcal O^\times.
 \tag{BX11}
\]
Suppose \(z:G_{\mathbb Q}\to k\) is a continuous
\(\gamma\)-cocycle with the following exact
\(\mathrm{UnrEv}\) condition:
at each finite inertia group, its restriction to
\(I\cap\ker\gamma\) is zero on a subgroup of
index prime to \(\ell\).
Then \(z\) is a coboundary.
\end{srlemma}
\begin{proof}
Reduction is injective on roots of unity of order prime
to \(\ell\), so \(\ker\chi=\ker\gamma\).
Let \(L\) be their fixed field.
Its degree over \(\mathbb Q\) is prime to \(\ell\).
On \(G_L\), the cocycle is an additive homomorphism.
On each \(I\cap G_L\) it has finite additive
\(\ell\)-group image; the prime-to-\(\ell\) index in the
hypothesis forces that image to be zero.
Thus this homomorphism is unramified at every finite
place.
At real places an element of order two has zero image
as well.
Class field theory makes it a homomorphism
\[
                         F:\operatorname{Cl}(L)\longrightarrow k .
 \tag{BX12}
\]
The cocycle identity gives
\(F(\sigma A)=\gamma(\sigma)F(A)\).

We recall the exact classical annihilator used on (BX12).
For conductor \(F_\chi\), put
\[
 \Theta=\frac1{F_\chi}
       \sum_{u\in(\mathbb Z/F_\chi)^\times}
                       u_0\,\sigma_u^{-1}.
\]
For every integer \(b\) prime to \(F_\chi\),
\((b-\sigma_b)\Theta\) is integral, and its projection
to \(\mathbb Z[\operatorname{Gal}(L/\mathbb Q)]\)
annihilates \(\operatorname{Cl}(L)\).
This is the integral Stickelberger annihilation theorem,
which applies to every abelian subfield, not merely the
full cyclotomic field.\footnote{For this precise abelian
subfield formulation see N.~Aoki,
\href{https://archive.mpim-bonn.mpg.de/id/eprint/743/1/preprint_1990_16.pdf}
{\emph{On the Stickelberger ideal of a composite field of
some quadratic fields}}, Theorem 1.1 and the restriction
definition immediately preceding it.
The theorem is the classical Stickelberger relation.}
Its integrality can also be seen immediately:
the coefficient at \(\sigma_u^{-1}\) is
\(\lfloor bu_0/F_\chi\rfloor\).

Choose \(b\) so that
\[
                         b-\bar\chi(b)\ne0\quad\text{in }k .
 \tag{BX13}
\]
If \(\ell\nmid F_\chi\), take \(b=\ell\).
If \(\ell\mid F_\chi\), condition
\(\gamma\ne\bar\epsilon_\ell\), cyclotomic
surjectivity and a unit representative modulo \(F_\chi\)
supply such a \(b\).
Evaluating the integral annihilator at \(\chi\) gives
\[
           (b-\chi(b))B_{1,\chi_{\rm prim}^{-1}} .
 \tag{BX14}
\]
It is a unit.
To include an imprimitive displayed level, use
\[
 B_{1,\chi^{-1}}
   =B_{1,\chi_{\rm prim}^{-1}}
            \prod_{\substack{p\mid N\\p\nmid F_\chi}}
                                  (1-\chi_{\rm prim}^{-1}(p)).
\]
The primitive Bernoulli number is integral because
(BX13) and integrality of (BX14) permit division by
the unit \(b-\chi(b)\).
All the Euler factors are integral.
The unit hypothesis in (BX11) therefore makes the
primitive Bernoulli number a unit.
This proves the assertion about (BX14).

Apply \(F\) to the annihilator acting on an arbitrary
ideal class.
Equivariance in (BX12) multiplies \(F(A)\) by the
nonzero reduction of (BX14).
Thus \(F=0\), so \(z|_{G_L}=0\).
The class of \(z\) descends to
\(H^1(\operatorname{Gal}(L/\mathbb Q),k(\gamma))=0\):
averaging is valid since the group order is prime to
\(\ell\).
Hence \(z\) is a coboundary, as asserted.
\end{proof}

\begin{srproposition}\label{prop:sr-actual-corx}
For all the \(\mathrm{XData}\) in the actual Case II
and large-conductor Gaussian constructions, the residual
cocycles in \(\mathrm{CorX}\), for both the indicated
character and its inverse, are coboundaries.
\end{srproposition}
\begin{proof}
In general \(\mathrm{XData}(\ell,d,\phi,\xi)\) is
exactly (BX1), together with the stated primitivity,
order and positivity conditions in
Lemma~\ref{lem:sr-dyadic-bernoulli-unit}.
That lemma also applies to \(\phi^{-1},\xi^{-1}\):
their conductors and orders are unchanged, and their
parity values at \(-1\) are unchanged.
Thus both required Bernoulli numbers are units.

The product character is nontrivial and of order
prime to \(\ell\).
It differs residually from the cyclotomic character.
Indeed choose a dyadic unit on which \(\xi=-1\),
put the odd-modulus residues equal to one, and also
put the residue modulo \(\ell\) equal to one.
The Chinese remainder theorem gives a global cyclotomic
element with \(\bar\chi=-1\) and
\(\bar\epsilon_\ell=1\).
Such a dyadic unit exists because the nontrivial
\(2\)-power character has \(-1\) in its image.
The same witness works for the inverse.
Lemma~\ref{lem:sr-bernoulli-unit-residual-unramified}
therefore proves exactly both clauses of
\(\mathrm{CorX}\), with the source's
\(\mathrm{UnrEv}\) condition retained.

We check that the actual size choices supply (BX1).
For Case II the required residual odd roots are
the characters
\(\bar\xi,\bar\Psi\bar\xi^{-1},
  \bar\Psi\bar\xi\), and their inverses.
Their odd-conductor primitive factors have conductors
\(f\mid\ell N\) and orders dividing
\(\ell^{d_0}-1\), with
\[
              d_0=\operatorname{ord}_4(\ell)\le2,\qquad
                       e=v_2(\ell^{d_0}-1)=w .
\]
Multiplication by the original dyadic character does
not lower the conductor of the primitive auxiliary
dyadic character, since \(n>c_2(\psi)\).
Their parity is odd: \(\Psi(c)=1\) while the
auxiliary dyadic character is odd.
The chosen bound
\[
                         2^n\ge\ell N4^w
\]
therefore gives \(f4^e\le2^n\).
This supplies the residual vanishing used by
\srcid{SymplecticGSp4.DataEx}, rather than
assuming its \(\mathrm{Raw.Good}\) conclusion.

For the Gaussian host the required corner character is
\[
                    \bar\Psi^{\,2}\eta_K
                         =\bar\epsilon_\ell^{\,2}\bar\psi^{\,2}\eta_4.
\]
The odd-prime components of the principal \(\psi\)
have order at most two, so they disappear after squaring.
Its odd primitive conductor is consequently at most
\(\ell\), not \(qd\ell\).
The dyadic factor \(\psi_2^2\eta_4\) is primitive of
conductor \(2^{c-1}\), since \(c\ge14\).
It is odd.
With \(d_0=2\), \(w=v_2(\ell^2-1)\), the large-branch
condition gives
\[
        2^{c-1}\ge R_1(\ell)/2
                  =8\ell^2 4^w\ge\ell4^w .
\]
Thus (BX1) holds here too.
The source's \(\mathrm{HostXpp}\) is the restriction
of this residual statement to the actual quotient.
More explicitly put
\(\kappa=(\eta_K,\bar\Psi)\) and
\(\gamma=(\bar\Psi^2\eta_K)^{-1}\).
The finite image of \(\kappa\) has order prime to
\(\ell\), and \(\ker\kappa\subset\ker\gamma\).
Away from \(\ell\), the actual host quotient makes
the pulled-back corner cocycle zero on
\(I_p\cap\ker\kappa\).
Its index in \(I_p\cap\ker\gamma\) divides
\(|\operatorname{im}\kappa|\), hence is prime to
\(\ell\).
At \(\ell\), the imposed full-local-group
coboundary is zero on \(I_\ell\cap\ker\gamma\).
These are exactly the \(\mathrm{UnrEv}\) conditions.
The proved residual lemma makes the pulled-back cocycle
a coboundary; surjectivity of the host quotient map
gives the same formula on the quotient itself.
This matches \srcid{LevelCohomology/Xpp.lean}
without assuming its conclusion.

All these estimates depend on the original conductor
and \(\ell\), with the already fixed exponents.
No additional field-degree or class-rank bound is introduced.
\end{proof}

% Frozen source:
% LevelCohomology/LxPartial, LxCore, LemmaX,
% ResidualBernoulliNonvanishing, ChrInvNeg, HcXvan, Xpp;
% SymplecticGSp4/DataEx and DirichletPrimeToLPartMulLPowerPart.
